\chapter{Evaluation and Validation}
\label{ch:eval}

\section{Baselines}

\subsection{Physics-Only Baseline}

The reduced-order physics model serves as a baseline, providing deterministic predictions without ML enhancement. This baseline represents the state-of-the-art for fast simulation in many rail applications.

% TODO: Report physics baseline metrics (MAE, RMSE, calibration)

\subsection{Non-Transformer Surrogate (Optional)}

For comparison, we also train a feedforward neural network surrogate that does not use transformer architecture. This baseline helps isolate the contribution of attention mechanisms.

% TODO: If implemented, report comparison metrics

\section{Metrics}

\subsection{Accuracy Metrics}

For deterministic outputs (energy, arrival time):
\begin{itemize}
    \item Mean Absolute Error (MAE): $\text{MAE} = \frac{1}{N} \sum_{i=1}^N |y_i - \hat{y}_i|$
    \item Root Mean Squared Error (RMSE): $\text{RMSE} = \sqrt{\frac{1}{N} \sum_{i=1}^N (y_i - \hat{y}_i)^2}$
    \item Mean Absolute Percentage Error (MAPE): $\text{MAPE} = \frac{100}{N} \sum_{i=1}^N \left|\frac{y_i - \hat{y}_i}{y_i}\right|$
\end{itemize}

% TODO: Report metrics in results table

\subsection{Calibration Metrics}

For probabilistic outputs:
\begin{itemize}
    \item Expected Calibration Error (ECE): Measures deviation from perfect calibration
    \item Reliability diagrams: Visualize predicted vs empirical quantile coverage
    \item Sharpness: Width of prediction intervals (tighter is better if calibrated)
\end{itemize}

% TODO: Add reliability diagram figure
% TODO: Report ECE values for each quantile level

\subsection{Constraint Violation Rates}

For safety-critical outputs (coupler forces), we measure:
\begin{itemize}
    \item False negative rate: Fraction of exceedances that were not predicted
    \item False positive rate: Fraction of predicted exceedances that did not occur
    \item Coverage at threshold: Fraction of true exceedances captured by $95\%$ prediction interval
\end{itemize}

% TODO: Report violation rates in table
% TODO: Discuss implications for safety margins

\section{Stress Tests}

\subsection{Steep Grades}

We evaluate performance on routes with grades exceeding \SI{2}{\percent}, where coupler forces are most critical. The model should maintain accuracy and calibration even under extreme conditions.

% TODO: Report metrics on steep grade subset
% TODO: Add figure showing coupler force predictions on steep grade

\subsection{Long and Heavy Consists}

Consists with 150+ cars or total mass exceeding \SI{10000}{\tonne} test the model's ability to handle force propagation through long chains. Attention mechanisms should enable modeling long-range dependencies.

% TODO: Report performance vs consist length
% TODO: Analyze attention patterns for long consists

\subsection{Low Adhesion Conditions}

Low adhesion scenarios (friction coefficient $< 0.15$) are safety-critical as they can lead to wheel slip and coupler overload. The model must provide conservative uncertainty estimates to avoid underestimating risk.

% TODO: Report calibration under low adhesion
% TODO: Compare predicted vs actual exceedance rates

\section{Ablation Studies}

\subsection{Without Physics Residual}

We train a variant that predicts outputs directly rather than as residuals to the physics layer. This ablation quantifies the benefit of physics integration.

% TODO: Report comparison: residual vs direct prediction
% TODO: Add table showing ablation results

\subsection{Without Factorized Attention}

We compare factorized attention to full attention (when computationally feasible) to assess the trade-off between expressiveness and efficiency.

% TODO: Report accuracy vs computational cost trade-off

\subsection{Without Calibration Term}

Removing the calibration loss term tests whether explicit calibration is necessary or if quantile regression alone provides sufficient calibration.

% TODO: Report ECE with/without calibration loss
% TODO: Add reliability diagrams for both variants

\section{Results Summary}

% TODO: Create comprehensive results table comparing all methods
% TODO: Include figures:
%   - Prediction accuracy plots (scatter, time series)
%   - Calibration diagrams
%   - Attention visualization
%   - Ablation comparison charts

Table~\ref{tab:main_results} summarizes the main evaluation results. Figure~\ref{fig:calibration} shows reliability diagrams for quantile predictions.

\begin{table}[h]
    \centering
    \caption{Main evaluation results comparing hybrid framework to baselines}
    \label{tab:main_results}
    \begin{tabular}{lccc}
        \toprule
        Method & Energy MAE (\si{\kilo\watt\hour}) & Arrival Time MAE (\si{\second}) & Coupler Force ECE \\
        \midrule
        Physics-only & TODO & TODO & TODO \\
        Hybrid (ours) & TODO & TODO & TODO \\
        \bottomrule
    \end{tabular}
\end{table}

\begin{figure}[h]
    \centering
    \includegraphics[width=0.8\textwidth]{figures/calibration_diagram.pdf}
    \caption{Reliability diagrams showing calibration of quantile predictions}
    \label{fig:calibration}
\end{figure}
